% ---
% id: receiver/rem:positivity
% kind: rem
% title: "Source of positivity of the click-pattern elements"
% summary: >-
%   The inclusion-exclusion formula is an alternating sum and nothing in its
%   form makes it positive. Positivity comes first, from commuting projective
%   measurements pulled back through an isometry, and the single-click case
%   also follows from monotonicity of the lift.
% prominence: standard
% status: accepted
% uses: [thm:click-povm, def:physmodel, lem:fock-monotone]
% assumes: []
% refs: []
% provenance: newly-derived
% lean:
% ---

\begin{remark}
\label{rem:positivity}
Formula \eqref{eq:incexc} is an alternating sum of positive operators and
nothing about its form makes it positive; the element of a non-empty outcome
$S$ has $2^{|S|}$ terms, half of them with a minus sign. Positivity is established \emph{before} the expansion, in
part~(i) of Theorem~\ref{thm:click-povm}, from the physical construction: $K$
commuting binary projective measurements on disjoint output mode sets, pulled
back through an isometry. Definition~\ref{def:physmodel} is set up for that
order of argument. A derivation that started from \eqref{eq:incexc} would
have to prove positivity of an alternating sum of Fock functors, which is both
harder and unnecessary. For a small independent check,
$\eqref{eq:Msc}\succeq0$ follows separately from
Lemma~\ref{lem:fock-monotone}, since $A_{[K]\setminus\{k\}}\preceq A_\Sigma$ gives
$\Fock(\id-A_{[K]\setminus\{k\}})\succeq\Fock(\id-A_\Sigma)$.
\end{remark}
